\subsection{P1 lấy độ lớn và tìm chuyển tiếp}

Kiến trúc Parser dùng chung cấu trúc \texttt{posit\_parser\_mul} cho cả ba toán hạng A, B và C.
Ở P1, dấu và hai cờ đặc biệt đến từ mã gốc. Độ lớn được tạo bằng
XOR với mặt nạ dấu và cộng bù, dùng đường carry của Artix-7.
Chuỗi XOR giữa hai bit kề nhau đánh dấu vị trí thay đổi từ run
đầu sang terminator. Vì vậy cùng một detector tìm chuyển tiếp
cho cả run 0 và run 1. Thêm sentinel 1 ở cuối giúp định nghĩa
được kết quả khi mã độ lớn không có terminator trong 31 bit.

\rtlfig{Hai tầng của \sig{posit\_parser\_mul}. Vector chuyển tiếp chỉ cấp LOD; seed để dịch giữ nguyên từ magnitude. Quire dùng \sig{posit\_parser} với cùng phép dò regime và trả fraction 26 bit.}{fig:parser-stages}{
\node[pin] (p) at (0,.42) {\sig{p[31:0]}};
\node[unit,text width=62mm] (mag) at (-3.05,-1.1) {\hw{Lấy độ lớn}\\$m=(p\oplus\{32\{s\}\})+s$};
\node[unit,text width=52mm] (exc) at (4.2,-1.1) {\hw{Dấu và đặc biệt}\\$s=p[31]$, $u=\lnot (|p[30:0])$\\$Z=\lnot s\land u$, $N=s\land u$};
\draw[wire] (p.south)--++(0,-.2) -| (mag.north);
\draw[wire] (p.south)--++(0,-.2) -| (exc.north);
\node[unit,text width=62mm] (xor) at (-3.05,-2.95) {\hw{Đánh dấu chuyển tiếp}\\$d=\{m[29:0]\oplus m[30:1],1'b1\}$};
\draw[wire] (mag)--node[tag,right]{$m[31:0]$}(xor);
\node[unit,text width=62mm] (lod) at (-3.05,-4.47) {\hw{LOD regime 31 bit}\\\sig{posit\_parser\_f\_regime\_lod}};
\draw[wire] (xor)--node[tag,right]{$d[30:0]$}(lod);
\node[ff,text width=125mm] (p1) at (0,-6.18) {\hw{P1} \quad \sig{cnt[4:0]} \enspace|\enspace
 $r=m[30]$ \enspace|\enspace \sig{seed[28:0]} \enspace|\enspace $s,Z,N$};
\draw[wire] (lod)--node[tag,right]{\sig{cnt[4:0]}}(p1.north -| lod);
\coordinate (seedport) at (.85,0);
\draw[wire] (mag.east) -- (.85,-1.1) -- (p1.north -| seedport)
    node[pos=.55,tag,right]{$m[30], m[28:0]$};
\draw[meta] (exc.south)--(p1.north -| exc) node[pos=.5,tag,right]{$s,Z,N$};
\node[unit,text width=50mm] (sh) at (-3.66,-7.84) {\hw{Dịch trái 29 bit}\\$w=\sig{seed}\ll\sig{cnt}$};
\node[unit,text width=52mm] (reg) at (3.66,-7.84) {\hw{Tái tạo regime signed}\\$r=1$: $k=\sig{cnt}$\\$r=0$: $k=-\sig{cnt}-1$};
\draw[wire] (p1.south -| sh)--(sh);
\draw[wire] (p1.south -| reg)--(reg);
\node[unit,text width=53mm] (slice) at (-3.66,-9.73) {\hw{Tách fraction}\\$f_H=w[25:14]$, $f_L=w[13:0]$\\\sig{tail\_cut} $=|f_L$};
\node[unit,text width=52mm] (scale) at (3.66,-9.73) {\hw{Ghép scale}\\$sf=8k+e$, $e=w[28:26]$\\chọn $k=0$ nếu $Z\lor N$};
\draw[wire] (sh)--(slice);
\draw[wire] (reg)--(scale);
\draw[wire] (sh.east)--(0,-7.84) |- (scale.west) node[pos=.49,tag,right]{$e$:3};
\node[ff,text width=125mm] (p2) at (0,-11.49) {\hw{P2} \quad \sig{sf[10:0]} \enspace|\enspace
 \sig{frac[11:0]}, \sig{frac\_lo[13:0]} \enspace|\enspace \sig{tail\_cut}, $s,Z,N$};
\draw[wire] (slice)--(p2.north -| slice);
\draw[wire] (scale)--(p2.north -| scale);
\draw[meta] (p1.east)--(6.91,-6.18)--node[tag,left]{$s,Z,N$}(6.91,-11.49)--(p2.east);
\node[pin] (ab) at (-3.66,-12.81) {A/B: $f_H$, \sig{tail\_cut}};
\node[pin] (c) at (3.66,-12.81) {C: \sig{\{frac,frac\_lo\}}};
\draw[wire] (p2.south -| ab)--(ab);
\draw[wire] (p2.south -| c)--(c);
}



P1 chốt chỉ số \texttt{cnt}, bit regime đầu, seed \texttt{mag[28:0]}
và dấu/cờ. Seed lấy từ độ lớn gốc. Vector XOR dùng riêng để dò
vị trí và không đi vào đường fraction. Đây là điểm quan trọng nếu
đối chiếu hình với RTL: dữ liệu dịch ở P2 là seed của magnitude,
không phải kết quả XOR. Với mẫu gần cực trị, sentinel giúp detector
có đầu vào khác 0 để chọn offset xác định.

\subsection{LOD/LZD: quy ước, cấu trúc và vị trí sử dụng}

Trong thiết kế phần cứng Posit \cite{jaiswal2018posit,murillo2020flopoco}, bộ đếm
bit~0 dẫn đầu (LZD) và bộ đếm bit~1 dẫn đầu (LOD) thường được sử dụng song song
để xác định độ dài chuỗi \textit{regime} (lần lượt tìm bit~1 hoặc bit~0 đầu
tiên sau chuỗi lặp) cũng như chuẩn hóa kết quả cộng/trừ. Trong RTL của đề tài,
khối \sig{lod} được chuẩn hóa theo quy ước tìm \emph{vị trí bit 1 đầu tiên} từ MSB.
Khi \sig{vld=1}, vị trí $K$ tính từ 0 cũng chính là số bit 0 dẫn đầu
(LZC). Với đầu vào toàn 0, bộ dò trả $K=0$, \sig{vld=0}.

Parser chuyển cả hai bài toán đếm run về cùng một phép dò:
XOR hai bit kề tạo bit 1 tại chỗ đổi cực tính đầu tiên, dù
run bắt đầu bằng 0 hay 1. Khi đó độ dài run là $K+1$.
Nếu không có chuyển tiếp, sentinel làm $K=30$ và
\sig{vld=1}. Cách này dùng một đường dò chung cho hai cực tính.

Hình~\ref{fig:lod-regime-tree} khai triển cây gộp radix-4 trong
Parser A/B/C và Parser của quire. Hình~\ref{fig:lod-sac-adder}
cho thấy hai cách cài đặt khác: SAC dùng cây nhị phân bốn cấp
cho fraction 12 bit; S4 và Parser tích NF dùng \sig{casez} 31 bit.
Tên file \sig{lod\_lzd\_core\_fixed2} và \sig{fixed4}
không biểu thị radix. Đường READ của quire có phép dò riêng trên
word 36 bit (Hình~\ref{fig:quire-context}).

\rtlfig{LOD tìm chuyển tiếp regime của Parser A/B/C và quire: ba cấp gộp 4:1. Nhánh đệm hằng được vẽ để thể hiện cấu trúc khai triển của RTL.}{fig:lod-regime-tree}{
\node[pin] (in) at (0,0) {$d[30:0]=\{m[29:0]\oplus m[30:1],1'b1\}$};
\node[unit,text width=115mm] (pad) at (0,-1.0) {\hw{Căn trái vector tìm kiếm}\\
  \sig{padded[63:0]} $=\{d[30:0],33'b0\}$};
\draw[wire] (in)--(pad);
\foreach \n/\x/\r/\s in {3/-5.4/63:48/16 bit dữ liệu,2/-1.8/47:32/15 bit + 0,1/1.8/31:16/16 bit 0,0/5.4/15:0/16 bit 0}{
  \node[pad,text width=27mm] (g\n) at (\x,-2.68) {\sig{[\r]}\\ \s};
  \draw[wire] (pad.south -| g\n)--(g\n);
  \node[unit,text width=27mm] (a\n) at (\x,-4.24) {\hw{Cấp 1}\\$4\times$ gộp 4:1};
  \draw[wire] (g\n)--(a\n);
  \node[unit,text width=27mm] (b\n) at (\x,-5.93) {\hw{Cấp 2}\\gộp 4:1};
  \draw[wire] (a\n)--node[tag,right]{$4\times(v,p_2)$}(b\n);
}
\node[unit,text width=133mm] (root) at (0,-7.64) {\hw{Cấp 3: chọn nhóm hợp lệ cao nhất}\\
   \sig{pos[5:0]} $=\{g_3[1:0],p_4[3:0]\}$};
\foreach \n in {0,1,2,3}{\draw[wire] (b\n)--node[tag,right]{$(v,p_4)$}(root.north -| b\n);}
\node[pin] (out) at (0,-8.91) {\sig{vld} $=v_{\mathrm{root}}$ \qquad \sig{K[4:0]} $=$ \sig{vld ? pos[4:0] : 0}};
\draw[wire] (root)--(out);
\node[explain,text width=133mm] at (0,-10.25) {\hw{Một nút gộp}\\
 $v=\bigvee_{i=0}^{3}v_i$. Ưu tiên nhóm $3,2,1,0$; mã vị trí nhóm lần lượt $00,01,10,11$.\\
 Ghép mã nhóm ở phía cao và vị trí con được chọn ở phía thấp.};
\node[pad,text width=133mm] at (0,-12.1) {\hw{Ví dụ từ MSB: cùng detector cho hai dạng terminator}\\
 \begin{tabular}{@{}lll@{}}
   Đếm bit 1:  \sig{11101} & $\to$ tìm bit 0 đầu tiên & \sig{XOR = 0011}, $K=2$ \\
   Đếm bit 0:  \sig{00010} & $\to$ tìm bit 1 đầu tiên & \sig{XOR = 0011}, $K=2$
 \end{tabular}\\[2pt]
 Không có chuyển tiếp: sentinel cho $K=30$, \sig{vld=1}.};
}


\rtlfig{Hai bộ dò còn lại trong datapath NF/F: cây nhị phân 12 bit cho SAC và mã hóa ưu tiên 31 bit cho adder và Parser tích NF.}{fig:lod-sac-adder}{
\node[head] at (-3.78,0) {\sig{lod\_lzd\_core\_fixed4}\\SAC: dò trên mask fraction};
\node[head] at (3.78,0) {\sig{lod\_lzd\_core\_fixed2}\\Adder S4 và giải mã tích NF};
\node[pin] (fx) at (-3.78,-.97) {\sig{fx[11:0]}};
\node[pin] (in) at (3.78,-.97) {\sig{in[30:0]}};
\node[unit,text width=60mm] (pad) at (-3.78,-2.02) {\hw{16 lá bit}\\\sig{data\_pad = \{fx,4'b0\}}\\mỗi lá: $(v,p)=(\text{bit},0)$};
\draw[wire] (fx)--(pad);
\foreach \l/\y/\n/\step in {1/-3.49/8/1,2/-4.77/4/2,3/-6.05/2/4,4/-7.33/1/8}{
  \node[unit,text width=60mm] (b\l) at (-3.78,\y) {\hw{Cấp \l: \n\ cặp gộp} \quad \sig{STEP}=$\step$};
}
\draw[wire] (pad)--(b1);
\draw[wire] (b1)--(b2);
\draw[wire] (b2)--(b3);
\draw[wire] (b3)--(b4);
\node[explain,text width=60mm] (merge) at (-3.78,-8.96) {\hw{Quy tắc tại một cặp}\\
 $v=v_H\lor v_L$;\quad $p=v_H?p_H:(p_L\mathbin{|}\mathrm{STEP})$};
\draw[wire] (b4)--node[tag,right]{$(v,p)$}(merge);
\node[pin] (out1) at (-3.78,-10.34) {\sig{vld}, \sig{K[3:0]} $=\sig{vld}?p:0$};
\draw[wire] (merge)--(out1);
\node[unit,text width=60mm] (v) at (3.78,-2.02) {\hw{OR 31 bit}\\\sig{vld = |in}};
\node[unit,text width=60mm] (pz) at (3.78,-4.72) {\sig{casez}:~\hw{mẫu từ MSB}\\[1pt]
 \sig{1???...} $\to K=0$\\
 \sig{01??...} $\to K=1$\\
 \sig{001?...} $\to K=2$\\
 \sig{...}\\
 \sig{00...001} $\to K=30$\\[1pt]
 mặc định $K=0$};
\draw[wire] (in.south)--(v.north);
\draw[wire] (in.east)--++(2.9,0)|-(pz.east);
\node[unit,text width=60mm] (out2) at (3.78,-7.25) {\hw{Hai đầu ra tổ hợp}\\\sig{vld}, \sig{K[4:0]}};
\draw[wire] (pz)--node[tag,right]{$K$}(out2);
\draw[wire] (v.west)--++(-.37,0)|-(out2.west) node[pos=.46,tag,right]{\sig{vld}};
\node[pad,text width=60mm] at (3.78,-8.96) {\hw{Cùng mã hóa, hai nguồn vào}\\
 S4: \sig{magnitude[30:0]}\\
 NF: vector XOR kề + sentinel};
\node[tag] at (3.78,-10.34) {Ví dụ: \sig{in=0} $\Rightarrow$ \sig{vld=0, K=0}};
}



\subsection{P2 căn payload và tạo hai phần fraction}

P2 vừa dịch seed sang trái, vừa tạo regime signed từ \texttt{cnt}
và polarity. Khi regime bắt đầu bằng 1, $k=\texttt{cnt}$; khi bắt
đầu bằng 0, $k=-\texttt{cnt}-1$. Ba bit trên của payload dịch là
exponent. Phần còn lại được tách thành $f_H[11:0]$ và $f_L[13:0]$.
P2 ghi cả hai phần và \texttt{tail\_cut} bằng OR của $f_L$.
Ở A/B, phần thấp chỉ đóng góp thông tin cắt. Ở C, hai phần được
ghép lại để bộ cộng nhận đủ độ chính xác khả dụng.

Ví dụ FT-05, \texttt{1c900000} giải mã $sf_A=-10$ và
$F_A=\texttt{0x480}$. \texttt{3c820000} giải mã $sf_B=-1$
và $F_B=\texttt{0x208}$. Hai mã cùng dấu dương và phần thấp đã
bỏ bằng 0. Bộ chọn OPS thấy cả hai fraction có hai bit 1, nên
chọn A làm X và B làm Y. Cách kiểm các số này là giải mã raw
độc lập, sau đó so dấu, tổng scale và magnitude ở trace core.

\subsection{Điều kiện advance qua P1/P2}

Nếu P2 rỗng hoặc downstream nhận dữ liệu, P2 có thể chốt payload
mới từ P1. P1 có thể nhận thêm khi P1 rỗng hoặc khi P2 tiến.
Thanh ghi valid tiến cùng với dữ liệu. Đối với top hiện hành,
ba Parser retire cùng sự kiện \texttt{pair\_retire}; vì vậy
dữ liệu C không được ghép với tích từ request khác. Zero và NaR
đi theo metadata đến wrapper. Fraction hoặc scale của một mã
đặc biệt chỉ được sử dụng theo nhánh xử lý đã chọn.
